\section{Clock triage}
\label{sec:clock}

\begin{theorem}[Clock triage on \(\Lam\)]
\label{thm:clock}
The certified H\'enon survivor has the following three clock properties.
\begin{enumerate}
  \item The unit-clock determinant is \(2\pi\iu\)-periodic in \(s\) and has
  \(O(T)\) zeros in every bounded real strip up to height \(T\).
  \item The stored generating-function action is not a strictly positive roof.
  \item The instability roof in \eqref{eq:roof} is positive, its periodic sums
  obey \(T_{p^r}=rT_p\), and it is non-lattice.
\end{enumerate}
\end{theorem}

\begin{proof}
For unit time, direct calculation from \eqref{eq:adjacency} gives
\begin{equation}
 \det(I-zA)=1-z-z^3-z^4.
 \label{eq:unit-z}
\end{equation}
Thus
\[
 D_{\rm unit}(s)=1-\e^{-s}-\e^{-3s}-\e^{-4s}
\]
is invariant under \(s\mapsto s+2\pi\iu\).  A nonzero exponential polynomial
has finitely many zeros in a compact fundamental rectangle.  Translating that
rectangle proves the \(O(T)\) bound.  In fact,
\[
 1-z-z^3-z^4=-(z^2+1)(z^2+z-1),
\]
so the vertical towers can be written explicitly.

For the action clock, put \(q=1/\sqrt6\) and consider
\[
 (q_0,q_1,q_2,q_3)=(-q,-q,q,q).
\]
It satisfies the H\'enon recurrence and is an admissible primitive period-four
orbit inside the certified rectangles.  For the generating function
\[
 S_6(q_i,q_{i+1})=q_iq_{i+1}-q_i+2q_i^3,
\]
each of
\(\sum_iq_iq_{i+1}\), \(\sum_iq_i\), and \(\sum_iq_i^3\) vanishes.  Its
total action is therefore zero, excluding strict positivity.

On \(\Lam\), the certificate gives \(\abs{x}\ge1/3\) and
\(\abs{m^u}\le1/2\).  Hence
\begin{equation}
 J^u(z)\ge12\abs{x}-r\abs{m^u(z)}
 \ge4-\frac{123}{224}=\frac{773}{224}>1.
 \label{eq:positive}
\end{equation}
The roof is positive.  Invariance of the unstable graph gives
\(u_{j+1}=a_ju_j\), where \(a_j=-12x_j-rm^u(z_j)\).  Around a periodic orbit,
\[
 \prod_{j=0}^{n_p-1}a_j=\Lambda_{u,p}.
\]
Taking logarithms of absolute values proves the periodic-sum identity and the
repetition law.

A H\"older roof is called \emph{lattice} if it is H\"older-cohomologous to
\(c k\) for some \(c>0\) and integer-valued \(k\); in particular, every
periodic sum then lies in \(c\mathbb Z\).  It remains to rule this out.  The
negative fixed point
\[
 q_*=-\frac{1+\sqrt7}{6}
\]
has unstable multiplier \(L_1>1\) with minimal polynomial
\begin{equation}
 P_1(X)=X^4-4X^3-22X^2-4X+1.
 \label{eq:p1}
\end{equation}
The explicit period-four orbit has trace \(578\) and multiplier
\[
 L_4=289+24\sqrt{145},
 \qquad
 P_4(X)=X^2-578X+1.
 \]
The four conjugates of \(L_1\) have distinct absolute values.  Consequently,
the four conjugates of \(L_1^m\) remain distinct for every \(m\ge1\), and
\([\Q(L_1^m):\Q]=4\).  Every \(L_4^n\) lies in
\(\Q(\sqrt{145})\) and has degree at most two.  An equality
\(\log L_4/\log L_1=m/n\in\Q\) would imply \(L_1^m=L_4^n\), a contradiction.
Two periodic lengths therefore have irrational ratio, which rules out a
common lattice.
The exact algebra and power-degree argument are expanded in
Appendix~\ref{app:algebra}.
\end{proof}

\begin{remark}[Scope of the theorem]
Non-lattice time removes the exact vertical periodicity forced by unit map
time.  It does not imply a meromorphic continuation, a functional equation, a
Riemann--von Mangoldt law, or an operator spectrum.  Those are separate
Route-A obligations.
\end{remark}
